% !TeX root = ../main.tex
\section{Converting logarithmic work to spatial stress}
\label{sec:logarithmic-derivative}

This appendix states the matrix derivative used in the finite-deformation
stress and Biot tensor. Let \(\mathbf C=\mathbf F^T\mathbf F\) and
\(\mathbf\varepsilon=\tfrac12\log\mathbf C\). If \(\mathbf T\)
is symmetric and conjugate to \(\mathbf\varepsilon\), then
\begin{align}
 \delta W&=\mathbf T:\delta\mathbf\varepsilon,
 \label{eq:logarithmic-conjugate-work}\\
 &=\left\{\mathbf F\left[\pder{\log\mathbf C}{\mathbf C}:
                         \mathbf T\right]\right\}:\delta\mathbf F.
 \label{eq:logarithmic-work-to-piola}
\end{align}
Thus the corresponding spatial Kirchhoff stress is
\begin{equation}
 \mathbf\tau=\mathbf F\left[\pder{\log\mathbf C}{\mathbf C}:
                         \mathbf T\right]\mathbf F^T.
 \label{eq:log-stress-pushforward}
\end{equation}
The same formula with \(\bar{\mathbf F}\) and
\(\mathbf T=\mathbb{C}_s:\bar{\mathbf\varepsilon}\) gives the intrinsic
mineral Kirchhoff stress \(\widehat{\mathbf\tau}_s\) in the true-deformation
frame. Equation \eqref{eq:rotated-mineral-kirchhoff} then maps it to the
mixture frame. The symbol \(\mathbf T\) is used only here to
state the derivative for an arbitrary symmetric conjugate stress.

In an orthonormal eigenbasis of \(\mathbf C\), with positive eigenvalues
\(\lambda_i\), the derivative acts componentwise as
\begin{equation}
 \left[\pder{\log\mathbf C}{\mathbf C}:\mathbf T\right]_{ij}
 =\begin{cases}
 \dfrac{\ln\lambda_i-\ln\lambda_j}{\lambda_i-\lambda_j}\,T_{ij},
                       &\lambda_i\ne\lambda_j,\\[5pt]
 T_{ij}/\lambda_i,      &\lambda_i=\lambda_j.
 \end{cases}
 \label{eq:log-frechet-spectral-form}
\end{equation}
The repeated-eigenvalue expression is the continuous limit, including
isotropic stretches. The coefficient is symmetric in \(i,j\), which makes
the derivative self-adjoint. The derivative maps \(\mathbf C\) to \(\mathbf I\).
Consequently,
\begin{align}
 \tr\mathbf\tau
 &=\mathbf C:\left[\pder{\log\mathbf C}{\mathbf C}:\mathbf T\right],
 \label{eq:logarithmic-stress-trace-expansion}\\
 &=\tr\mathbf T.
 \label{eq:logarithmic-stress-trace}
\end{align}
For \(\mathbf T=\mathbf I\), it maps the logarithmic stress to the same
spatial identity. In general the transformation retains the different
principal directions of strain and stress; it cannot be replaced by a
rotation of \(\mathbf T\).
